% Part: first-order-logic
% Chapter: axiomatic-deduction
% Section: proof-theoretic-notions

\documentclass[../../../include/open-logic-section]{subfiles}

\begin{document}

\iftag{FOL}
      {\olfileid{fol}{axd}{ptn}}
      {\olfileid{pl}{axd}{ptn}}

\olsection{প্রমাণতাত্ত্বিক ধারণা}

\begin{explain}
যেভাবে আমরা কয়েকটি গুরুত্বপূর্ণ অর্থতাত্ত্বিক ধারণা
(\iftag{FOL}{বৈধতা}{সর্বতঃসত্যতা}, অনুসিদ্ধান্ত, পরিতৃপ্তিযোগ্যতা)
সংজ্ঞায়িত করেছি, এবার তেমনি তাদের অনুরূপ \emph{প্রমাণতাত্ত্বিক
ধারণা} সংজ্ঞায়িত করি। এগুলি কোনো !!{structure}s-তে !!{sentence}s-এর
পরিতৃপ্তির সাহায্যে নয়, বরং কয়েকটি সূত্রের !!{derivability} বা
!!{nonderivability}-র সাহায্যে সংজ্ঞায়িত। এই ধারণাগুলি যে পরস্পরের
সঙ্গে মিলে যায়, তা এক গুরুত্বপূর্ণ আবিষ্কার। তাদের মিলই
\emph{বিশুদ্ধতা} ও \emph{পূর্ণতা উপপাদ্য}-র বক্তব্য।
\end{explain}

\begin{defn}[!!^{derivability}]
কোনো !!^a{formula}~$!A$ $\Gamma$ \emph{থেকে !!{derivable}}, অর্থাৎ
$\Gamma \Proves !A$, যদি $\Gamma$ থেকে এমন কোনো !!a{derivation}
থাকে যার শেষ সূত্র~$!A$।
\end{defn}

\begin{defn}[উপপাদ্য]
কোনো !!^a{formula}~$!A$ একটি \emph{উপপাদ্য}, যদি শূন্য সেট থেকে
$!A$-এর কোনো !!a{derivation} থাকে। $!A$ উপপাদ্য হলে $\Proves !A$
এবং না হলে $\Proves/ !A$ লিখি।
\end{defn}

\begin{defn}[সঙ্গতি]
!!{formula}s-এর কোনো সেট~$\Gamma$ ঠিক তখনই \emph{সঙ্গত}, যখন
$\Gamma\Proves/ \lfalse$; নইলে সেটি \emph{অসঙ্গত}।
\end{defn}

\begin{prop}[প্রতিবিম্ব ধর্ম]
\ollabel{prop:reflexivity}
$!A \in \Gamma$ হলে $\Gamma \Proves !A$।
\end{prop}

\begin{proof}
  শুধু !!{formula}~$!A$-ই $\Gamma$ থেকে~$!A$-এর কোনো !!a{derivation}।
\end{proof}

\begin{prop}[একঘেয়েতা]
\ollabel{prop:monotonicity}
$\Gamma \subseteq \Delta$ এবং $\Gamma \Proves !A$ হলে $\Delta
\Proves !A$।
\end{prop}

\begin{proof}
$\Gamma$ থেকে~$!A$-এর যেকোনো !!{derivation}, $\Delta$
থেকে~$!A$-এরও কোনো !!a{derivation}।
\end{proof}

\begin{prop}[পরিযায়িতা]
\ollabel{prop:transitivity}
$\Gamma \Proves !A$ এবং $\{!A\} \cup \Delta \Proves
!B$ হলে $\Gamma \cup \Delta \Proves !B$।
\end{prop}

\begin{proof}
  ধরা যাক, $\{!A\} \cup \Delta \Proves !B$। তাহলে $\{!A\} \cup
  \Delta$ থেকে কোনো !!a{derivation} $!B_1$, \dots, $!B_l = !B$ আছে।
  ওই !!{derivation}-এর কয়েকটি ধাপ সঠিক হবে এমন কোনো বিধির জন্য, যা
  আগের কোনো পংক্তি~$!B_i = !A$-র উল্লেখ করে। অনুমান অনুযায়ী,
  $\Gamma$ থেকে~$!A$-এর কোনো !!a{derivation} আছে, অর্থাৎ এমন কোনো
  !!a{derivation}~$!A_1$, \dots, $!A_k = !A$, যেখানে প্রত্যেক $!A_i$
  হয় স্বতঃসিদ্ধ, নয়তো~$\Gamma$-র একটি !!a{element}, নয়তো কোনো
  অনুমান-বিধি অনুসারে সঠিক। এবার অনুক্রমটি বিবেচনা করো:
  \[
  !A_1, \dots, !A_k = !A, !B_1, \dots, !B_l = !B.
  \]
  এটি $\Gamma \cup \Delta$ থেকে~$!B$-এর একটি সঠিক !!{derivation},
  কারণ প্রত্যেক $!B_i = !A$ এখন $!A_k = !A$-র মতো একই ভিত্তিতে
  সমর্থিত।
  % BN-SRC-058 TARGET-CORRECTION-BEGIN
  \par\smallskip\noindent\textnormal{\small\textbf{উৎস-সংশোধন (BN-SRC-058):} হিমায়িত ইংরেজি প্রমাণের শেষ বাক্যে সূত্র-চলকটির একবারের লেখায় উপসর্গচিহ্ন বাদ পড়ে $B_i = !A$ হয়েছে, যদিও এই প্রমাণের সর্বত্র সূত্রটি $!B_i$। বাংলা পাঠে $!B_i = !A$ পুনরুদ্ধার করা হয়েছে।} % BN-SRC-058 TARGET-CORRECTION-END
  % BN-SRC-070 TARGET-CORRECTION-BEGIN
  \par\smallskip\noindent\textnormal{\small\textbf{উৎস-সংশোধন (BN-SRC-070):} হিমায়িত ইংরেজি প্রমাণটি বলে, প্রতিটি পুনরাবৃত্ত $!B_i = !A$-কে সেই একই “বিধি” সমর্থন করে যা $!A_k = !A$-কে সমর্থন করে। কিন্তু $!A_k$ স্বতঃসিদ্ধ বা~$\Gamma$-র উপাদানও হতে পারে। বাংলা পাঠে তাই “একই ভিত্তিতে সমর্থিত” বলা হয়েছে: স্বতঃসিদ্ধ হলে স্বতঃসিদ্ধ হিসেবে, অনুমিতি হলে $\Gamma \subseteq \Gamma \cup \Delta$ অনুসারে, আর অনুমান-বিধি দিয়ে এলে আগের একই পংক্তিগুলি থেকে।} % BN-SRC-070 TARGET-CORRECTION-END
\end{proof}

লক্ষ করো, এর অর্থ বিশেষত এই: $\Gamma \Proves !A$ এবং $!A \Proves
!B$ হলে $\Gamma \Proves !B$। এর থেকে আরও পাওয়া যায়: $!A_1,
\dots, !A_n \Proves !B$ এবং প্রত্যেক~$i$-এর জন্য $\Gamma \Proves
!A_i$ হলে $\Gamma \Proves !B$।

\begin{prop}
\ollabel{prop:incons}
$\Gamma$ ঠিক তখনই অসঙ্গত, যখন প্রত্যেক~$!A$-এর জন্য $\Gamma \Proves !A$।
\end{prop}

\begin{proof}
অনুশীলনী।
\end{proof}

\tagprob{FOL}
\begin{prob}
\olref[fol][axd][ptn]{prop:incons} প্রমাণ করো।
\end{prob}
\tagendprob

\tagprob{notFOL}
\begin{prob}
\olref[pl][axd][ptn]{prop:incons} প্রমাণ করো।
\end{prob}
\tagendprob

\begin{prop}[সংহতি]
\ollabel{prop:proves-compact}
  \begin{enumerate}
  \item $\Gamma \Proves !A$ হলে এমন একটি সসীম উপসেট $\Gamma_0
    \subseteq \Gamma$ আছে যেন $\Gamma_0 \Proves !A$।
  \item $\Gamma$-র প্রত্যেক সসীম উপসেট সঙ্গত হলে $\Gamma$ সঙ্গত।
  \end{enumerate}
\end{prop}

\begin{proof}
  \begin{enumerate}
    \item $\Gamma \Proves !A$ হলে !!{formula}s-এর এমন একটি সসীম
      অনুক্রম $!A_1$, \dots,~$!A_n$ আছে যেন $!A \ident !A_n$ এবং
      প্রত্যেক $!A_i$ হয় কোনো যৌক্তিক স্বতঃসিদ্ধ, নয়তো~$\Gamma$-র
      কোনো !!a{element}, নয়তো আগের !!{formula}s থেকে অনুমোদিত কোনো
      অনুমান-বিধি অনুসারে অনুসৃত। যেসব $!A_i$~$\Gamma$-তে আছে,
      সেগুলিকে নিয়ে~$\Gamma_0$ গঠন করি। তাহলে একই !!{derivation}
      $\Gamma_0$ থেকেও কোনো !!a{derivation}; তাই $\Gamma_0 \Proves !A$।
      % BN-SRC-059 TARGET-CORRECTION-BEGIN
      \par\smallskip\noindent\textnormal{\small\textbf{উৎস-সংশোধন (BN-SRC-059):} হিমায়িত ইংরেজি সংহতি-প্রমাণটি প্রতিটি অনুসৃত ধাপকে কেবল মোডাস পোনেন্স থেকে আসা বলে, যদিও এই অধ্যায়ের FOL শাখায় পরিমাণসূচকের অনুমান-বিধিও বৈধ। বাংলা পাঠে “অনুমোদিত কোনো অনুমান-বিধি” বলা হয়েছে। সসীম অনুক্রমে ব্যবহৃত প্রত্যেক বিধির সসীম পূর্বধারণা থাকে এবং প্রারম্ভিক সেট ছোট করলে নতুন-ধ্রুবকের শর্ত বজায় থাকে, তাই একই সসীম-সমর্থন যুক্তি উভয় বিধির জন্য চলে।} % BN-SRC-059 TARGET-CORRECTION-END
    \item এটি (1)-এর বিশেষ ক্ষেত্র $!A \ident \lfalse$-এর
      বিপরীত-প্রতিজ্ঞা।
\end{enumerate}
\end{proof}

\end{document}
